\chapter{1939:Adolf Butenandt and Leopold Ruzicka}

\label{chap:chemistry1939}

The Nobel Prize in Chemistry for the year 1939 was awarded jointly to Adolf Butenandt and Leopold Ruzicka.

\textbf{Motivation:}

for his work on sex hormones (one half, Adolf Butenandt) and for his work on polymethylenes and higher terpenes (the other half, Leopold Ruzicka)

\section{Adolf Butenandt}

\subsection*{Early Life and Education}

Adolf Friedrich Johann Butenandt was born on 1903-03-24 in Lehe, Germany (now part of Bremerhaven).
He was a German biochemist who shared the 1939 Nobel Prize in Chemistry with Leopold Ruzicka for his work on sex hormones.

\subsection*{Career and Major Research}

Butenandt studied chemistry in Marburg and Göttingen and received his doctorate at the University of Göttingen in 1927 under Adolf Windaus.
Working in Windaus's laboratory, he turned to the sex hormones and in 1929 succeeded in crystallizing the female sex hormone estrone from the urine of pregnant women, the first sex hormone to be obtained in pure form.
He isolated the male sex hormone androsterone from urine and established its structure in 1934, showing that it is closely related to the female hormones.
He was professor at the Technische Hochschule in Danzig from 1933 to 1936.
In 1936 he became director of the Kaiser Wilhelm Institute for Biochemistry in Berlin-Dahlem, which moved to Tübingen in 1944 and later, as the Max Planck Institute for Biochemistry, to Munich.
Later he studied the insect moulting hormone ecdysone and the silkworm attractant bombykol.

\subsection*{Nobel Prize-winning Work}

The work for which Adolf Butenandt was awarded the Nobel Prize in Chemistry in 1939 was described by the Nobel Committee as: "for his work on sex hormones".

Butenandt isolated the female sex hormone estrone in pure crystalline form from pregnancy urine and later obtained the male sex hormone androsterone, establishing the structures of both and showing that the sex hormones form a single chemical family.

\subsection*{Significance and Impact}

Butenandt's isolation of the sex hormones made them available for study and for medical use, and his structural work revealed that the male and female hormones are closely related members of the steroid family.

\subsection*{Later Career and Honors}

In 1939 Butenandt was awarded the Nobel Prize in Chemistry but was obliged by the German government to decline it; he received the prize formally in 1949.
After the war he became one of the most influential figures in German science and served as president of the Max Planck Society from 1960 to 1972.
He died on 1995-01-18 in Munich, Germany.

\subsection*{Legacy}

The legacy of Adolf Butenandt endures through the lasting impact of his scientific contributions.
His isolation of the sex hormones laid the foundation of steroid endocrinology and of the pharmaceutical production of hormonal drugs.

\subsection*{Selected Publications}

\begin{itemize}

\item Butenandt, A. (1929). "Über 'Progynon', ein krystallisiertes weibliches Sexualhormon." Naturwissenschaften, 17, 879.

\item Butenandt, A., \& Tscherning, K. (1934). "Über Androsteron, ein krystallisiertes männliches Sexualhormon."

\end{itemize}

\clearpage

\section{Leopold Ruzicka}

\subsection*{Early Life and Education}

Leopold Ruzicka was born on 1887-09-13 in Vukovar, Croatia, then part of Austria-Hungary.
He was a Croatian-born Swiss organic chemist who shared the 1939 Nobel Prize in Chemistry with Adolf Butenandt for his work on sex hormones.

\subsection*{Career and Major Research}

Ruzicka studied at the Technische Hochschule in Karlsruhe, where he received his doctorate in 1910 under Hermann Staudinger, working on the chemistry of the ketenes.
He followed Staudinger to the Eidgenössische Technische Hochschule in Zürich, where he habilitated and took Swiss citizenship.
He was professor at the University of Utrecht from 1926 to 1929 and then professor of organic chemistry at the ETH in Zürich, where he built a large research school.
He established the constitution of the natural musk compounds muscone and civetone, showing that they contain large carbon rings, and he synthesized muscone, proving that rings of thirteen and seventeen carbon atoms are stable despite earlier claims to the contrary.
His investigations of the higher terpenes established the principles by which these substances are built from isoprene units.
With his co-workers he synthesized the male sex hormone testosterone from cholesterol in 1934-1935, applying his knowledge of terpenoid and steroid chemistry.

\subsection*{Nobel Prize-winning Work}

The work for which Leopold Ruzicka was awarded the Nobel Prize in Chemistry in 1939 was cited by the committee for his investigations on polymethylenes and higher terpenes.
The citation preserved for this chapter reads: "for his work on polymethylenes and higher terpenes".

Ruzicka's demonstration that natural products contain large carbon rings, his synthesis of muscone, and his investigations of the higher terpenes formed the body of work for which he was honoured.
Its culmination in the synthesis of the male sex hormone testosterone connected this research directly to the chemistry of the sex hormones.

\subsection*{Significance and Impact}

Ruzicka's work overturned a widely held belief about the stability of large carbon rings and opened the field of macrocyclic chemistry, which became important for the synthesis of perfumes and pharmaceuticals.
His synthesis of testosterone established a practical route to the steroid hormones.

\subsection*{Later Career and Honors}

Ruzicka remained at the ETH in Zürich until his retirement in 1957, and his school trained many organic chemists.
He held honorary doctorates from several universities and was elected a foreign member of the Royal Society.
He died on 1976-09-26 in Mammern, Switzerland.

\subsection*{Legacy}

The legacy of Leopold Ruzicka endures through the lasting impact of his scientific contributions.
Macrocyclic chemistry and the chemistry of the terpenes, fields he founded, continue to be built upon by organic chemists throughout the world.

\subsection*{Selected Publications}

\begin{itemize}

\item Ruzicka, L. (1926). "Zur Kenntnis des Kohlenstoffringes. Über die Konstitution des Muscons." Helvetica Chimica Acta.

\item Ruzicka, L., \& Wettstein, A. (1935). "Über die künstliche Herstellung des männlichen Sexualhormons Testosteron." Helvetica Chimica Acta.

\end{itemize}

\clearpage

\section*{Chapter Summary}

The 1939 prize recognized Adolf Butenandt for his work on sex hormones and Leopold Ruzicka for his work on polymethylenes and higher terpenes. Their investigations established the steroid and terpene chemistry of biologically active compounds.
